In the Gospel today the noise of the sea throws whole nations into confusion. Jesus tells the disciples that people will grow faint thinking about what is going to happen to the world. They already suffer from something they can only imagine. Then Jesus speaks of someone they will see. The Son of Man will come in a cloud with great power and majesty. When all this begins, he tells his disciples, they can raise their heads, “because your redemption is at hand.”

Jesus says this while they can still hear his voice and see his face. They have a chance now to learn from him what they will need to remember then. The man speaking to them is the one who will come in the cloud. He wants his friends to know something about their future that they could never learn from looking at the sea. He is coming to set them free.

We began Advent today by asking him to come. We spoke quite plainly about what we need him to free us from. In the Collect the priest asked the Lord to come to the rescue of people whose own sins have put them in danger. We said Amen to that. When someone has done us an injury, it can be a relief to think of a judge who knows the whole truth. But when the truth includes the harm we ourselves have done, we may well wish we had a little longer before we had to stand before him. Yet today we have asked him to come to us precisely because he knows what we have done and has the power to release us from it.

Saint Gregory preached this Gospel to people whose neighbors had died in a storm. Some of the dead had been making plans for the next day. Gregory asks them to picture a traveler worn out by the road who still wants the journey to go on. His hearers had good reason to want relief. Gregory asks them to consider how much they want to see the Lord at the end of the road.

There are days when any of us might want to have our troubles over and done with. Jesus speaks about something we can look forward to even when those troubles are only beginning. He says that we shall see him. He took from Mary the flesh in which he could give his life for us. When he speaks of our redemption, he knows what it will cost him to bring us home. Even someone who has long been afraid to face the Lord has reason to want that meeting. Before the Gospel we sang the words of people asking God to show them his mercy. We have seen what that mercy cost him on the cross. We can begin to want his company even while we are still afraid to tell him the whole truth about ourselves.

Saint Paul is speaking to people who already believe this. They can remember when they first believed, and Paul tells them that their salvation is nearer now. Perhaps you can remember a time when it came more readily to you to do what Christ asked. There may be things you once took pains to avoid that you have since grown quite easy about. Saint John Chrysostom asks Christians whose first eagerness has cooled to think of runners in a race. When the runners come within sight of the prize, they make a greater effort. Paul tells people who need waking up that what is coming nearer is their salvation. He still expects them to be able to get up and walk in the light.

Among the things he calls deeds of darkness are contention and envy. We can quarrel in broad daylight, with other people there who know exactly what was said. What can remain hidden for a long time is the pleasure one of us takes in keeping the quarrel alive. Suppose someone has spoken harshly to you. You tell a friend the words he used. Every word you report may be quite accurate. Yet you can tell the story in such a way that the person who hurt you will find it hard to be welcome in that friend's house again. The two of them will meet again, and you know it. Perhaps you have been looking forward to the reception he will get.

Just before Paul tells the Romans to wake from sleep, he has spoken to them about a debt they owe their neighbor. They owe him love. That is still true of the man who spoke harshly. In the Gradual we asked the Lord to teach us his paths. We are asking to be taught by someone who sees both the injury we suffered and the injury we are willing to inflict. If you have given a friend reason to despise someone you quarreled with, you can go back to that friend and tell the part of the story you left out about your own conduct. The other person may never learn that you did this. He may simply find it a little easier to be received in a place where you once made it hard for him.

In a few moments the priest will pray over the gifts at the altar. He will ask the Father to let these holy gifts cleanse us, so that we may come to him with purer hearts. Even when you have corrected the story, you may still very much want that man to feel what you felt. At Communion we sing about the goodness God will give and the fruit our earth will bear. Christ gives himself to us here. He can change even what it gives us pleasure to do. You might live to see that man warmly welcomed among your friends and find that you are glad he is there.

After Communion we will ask for mercy here in God's house as we prepare to celebrate the birth of his Son. Christmas is a day we can mark on the calendar. Jesus gives his friends a different kind of assurance about the day they will see him in majesty. He speaks to them about the fig tree and all the trees on which they have seen the first fruit. They recognize the approach of summer because they have seen fruit appear on those branches year after year. Jesus tells them that when they see the signs he has described, they will know that God's kingdom is near. They have never seen the day he is speaking of. For that they have his word, and he tells them it will stand even when heaven and earth have passed away. You can rely on that word when you have had to admit something about yourself that makes you ashamed. The Lord whose help you need today is the one you may hope to see face to face when he comes.
